\chapter{Conclusions and Future Work}
\label{chap:conclusions-future-work}

\section{Summary of the Study}
\label{sec:conclusions-future-work:summary-of-the-study}
% Section ID: C12.1.
% Evidence map: see thesis-writing/planning/chapter_evidence_map.md.

This thesis constructed and characterized two observational causal-analysis testbeds, instantiated separately from source-recorded MIMIC-III and PhysioNet 2012 measurements, covariates, care-process traces, and in-hospital mortality.  It did not simulate patient records, exposure-assignment mechanisms, or outcomes.  Instead, it made representation-defined proxy constructs, rules, cohorts, preprocessing, aggregation, empirical support, DAGs, adjustment assumptions, estimator interfaces, and provenance explicit.

Because neither resource contains a known causal answer key, characterization used converging, non-equivalent evidence rather than error against causal truth.  Integration, four-model proxy prediction, five-source aggregation, matching, three-estimator comparison, diagnostics, and evidence tracking made the resources inspectable.  Their proxy states remain analytical constructs rather than verified clinical conditions, and every adjusted result remains conditional on measurement, graph, intervention-definition, temporal-ordering, support, confounding, and modeling assumptions.

\section{Main Findings and Contributions}
\label{sec:conclusions-future-work:main-findings-contributions}

The strongest contribution is a representation-centered framework for constructing observational causal-analysis testbeds, together with two concrete ICU instantiations whose assumptions and provenance are exposed.  A shared proxy-state interface traces constructs from one rule-derived source through four predicted voters to separate dataset-specific exposure analyses, while causal and split-aware contracts keep identifiers, cohorts, and exports distinguishable.  Integration remains an enabling contribution: it makes the representation, estimator interfaces, and evidence boundaries auditable, but it is not itself the principal scientific object.  Checked sources and manifests improve numerical traceability without establishing clean-checkout reproduction.

The empirical predictive comparison covered STraTS, GRU, GRU-D, and TCN.  STraTS led the archived proxy-label test metrics in MIMIC-III, whereas GRU-D led them in PhysioNet (Chapter~\ref{chap:results}, Section~\ref{sec:results:predictive-performance}).  The leading model was therefore resource dependent; recoverability did not establish a universal ordering or clinical validity.  The separately admitted MIMIC proxy-vector result retained mortality-relevant predictive information, while the analogous PhysioNet number remained withheld because its lineage was unresolved.

In the primary original cohorts, CausalForestDML means were positive for all nine MIMIC-III exposures and nine of ten PhysioNet exposures; PhysioNet shock was negative.  CausalForestDML and LinearDML shared direction in all 19 dataset--exposure comparisons, and exploratory CausalPFN agreed with both in 18 of 19.  Within-resource rank correlations and pairwise RMSE showed bounded ordering and magnitude disagreement, while the nine-construct cross-resource ranking was only moderately correlated.  PhysioNet shock remained the central exception across negative DML, slightly positive PFN, positive matching, and a noncommensurate external comparison.

The clinical comparison retained both broadly compatible contexts and material discrepancies.  At the heuristic $\lvert z\rvert>2$ disruption threshold, 36/36 MIMIC and 34/40 PhysioNet DML comparisons passed; this is not a p-value or formal randomization test.  Matching failures, support warnings, outcome-downsampling changes, and heterogeneous sensitivity provenance prevented convergence from becoming closure.  Agreement shows that a pattern is not unique to one fitted estimator under one resource; it does not establish that the representation, graph, estimand, or estimate is correct.  The resources therefore complement controlled answer-key benchmarks while remaining unable to measure causal error against truth.

\section{Limitations of the Conclusions}
\label{sec:conclusions-future-work:limitations-of-the-conclusions}

Construct validity is the first constraint: LLM-assisted, project-selected proxy rules lack chart adjudication; dataset ontologies and timing differ; prediction can add error; and mixed-source voting can preserve shared rule and measurement error.  Clinical review, reference-standard evaluation, complete design-to-code decisions, and exact support for several rule families remain incomplete.

Identification and translation remain equally limited.  Illness-state proxies are not necessarily well-defined interventions; the project DAGs are assumptions; temporal ordering, overlap, uncertainty, and unmeasured confounding remain unresolved; and diagnostics do not establish causal validity.  The two datasets do not justify pooling or transport claims, and no prospective or external clinical validation, fairness or safety study, clinician-only comparison, alternative-LLM replication, or prompt-sensitivity analysis was performed.  No clinical or causal truth, treatment recommendation, or deployment readiness follows.  Ethics and governance wording and the data, configuration, split, checkpoint, source-version, environment, and archive records needed for full reproduction also remain incomplete.

\section{Future Work}
\label{sec:conclusions-future-work:future-work}

The first priority is clinical and construct validation.  Future work should freeze the proxy rules and conduct blinded multi-clinician review of every proxy definition and graph edge, with reported inter-reviewer agreement.  Rule-derived and predicted states should be compared with chart-adjudicated reference labels, including dataset-specific assessment of timing, missingness, thresholds, and differences among rule-derived labels, individual predicted labels, and mixed-source aggregated proxy labels.  Controlled clinician-only, LLM-assisted, and hybrid construction should be compared, together with alternative LLM families, prompt perturbations, model-version changes, and decoding choices.

The second priority is a stronger observational and measurement design.  A future study should specify a target-trial-aligned question with eligibility, time zero, intervention strategies, follow-up, outcome, estimand, and an explicit plan for time-varying confounding.  It should add prespecified overlap and balance diagnostics, support-aware estimands, uncertainty and multiplicity procedures, and richer sensitivity analyses.  Rule ablation and threshold-sensitivity experiments, formal label-noise modeling, and explicit measurement-error analyses should quantify how proxy construction, prediction, and aggregation affect downstream estimates.  DAG review should proceed edge by edge and include alternative-graph and adjustment-set analyses rather than treating one encoded graph as fixed truth.

The third priority is reproducibility and external evaluation.  Signed per-run manifests should link raw-data access records, processed-artifact hashes, split identifiers, checkpoints, commands, resolved configurations, software environments, source versions, and archived outputs.  LLM design manifests should record complete prompts, system and decoding settings, model versions, follow-up interactions, reviewer decisions, and prompt-output-to-code mappings.  Frozen definitions should then be evaluated on temporal, external, prospective, and additional held-out observational testbeds, including clinically justified subgroups.  Fairness, governance, privacy, safety, workflow, and human-factors evaluation must precede any deployment consideration.  Although its primary method source is now cited, CausalPFN should remain exploratory until the exact producing implementation version and checkpoint are verified and an appropriate diagnostic and uncertainty plan is available.  No clinical deployment should occur before construct, causal-design, external, and prospective validation and authoritative ethics and data-governance review.

\section{Closing Perspective}
\label{sec:conclusions-future-work:closing-perspective}

An LLM cannot replace a clinician or causal expert on the evidence presented here, and integration does not remove the need for careful observational design, diagnostics, and reproducible provenance.  Transparent representation and provenance make observational testbeds inspectable.  Their value lies not in supplying causal truth, but in exposing estimator stability, disagreement, representation-dependent support limitations, and sensitivity while retaining source-recorded measurements, missingness patterns, care-process traces, covariates, and outcomes.
